\section{Introduction}

Spontaneous eye blinks are simultaneously a nuisance source in scalp EEG and a potentially informative physiological signal. Blink-related EEG measures have been linked to cognitive load during driving, fatigue-related monitoring, and abnormal clinical states, which means that reliable blink extraction affects both preprocessing and interpretation.\supercite{alyan2023scirep,tarafder2022,wang2018,kimura2017,mentzel2017} \blinker{} was introduced to automate the extraction of ocular indices from EEG, electrooculography and independent-component time series, making large-scale blink analysis more tractable than manual inspection.\supercite{kleifges2017}

The practical limitation is that \blinker{} remains a MATLAB workflow. That matters because contemporary open-source EEG analysis increasingly centers on Python-native tools such as MNE-Python and the wider MNE ecosystem.\supercite{gramfort2013,appelhoff2019} A Python reimplementation is only useful, however, if it preserves the reference behavior closely enough that established \blinker-based analyses can be migrated without silently changing event tables. Manual or semi-manual alternatives do not remove this problem, because annotation itself is labor-intensive and can be error-prone even when nominal agreement is high.\supercite{jansen2020}

This requirement shifts the scientific question away from detector novelty and toward implementation fidelity. Small divergences in sample indexing, terminal-edge handling, or floating-point gating can materially change blink boundaries and therefore alter downstream ocular indices. For software intended to replace a reference implementation, approximate similarity is weaker evidence than exact parity on heterogeneous public data.

Here we present \tool, a Python implementation of \blinker{} designed for use within MNE-Python workflows, and evaluate it against the MATLAB reference across two public EEG resources that cover distinct recording contexts: motor-imagery recordings from the Murat 2018 dataset\supercite{kaya2018} and fatigue-driving recordings from the public SEED-VLA and SEED-VRW datasets.\supercite{luo2024} Our aim is deliberately narrow and testable: to determine whether the Python implementation reproduces MATLAB \blinker{} outputs exactly after full-dataset reruns, and to document the edge cases that had to be resolved to reach that result.
